\documentclass[letterpaper,11pt]{article}
\usepackage{graphicx}
\usepackage{geometry}
\linespread{1.2}                        
\setlength{\oddsidemargin}{0.0in}            
\setlength{\textwidth}{6.5in}    
\usepackage[utf8]{inputenc} % Asegúrate de usar utf8
\usepackage[T1]{fontenc} % Asegúrate de usar T1 font encoding      
\setlength{\topmargin}{-.6in}            
\setlength{\textheight}{9in}
\usepackage{latexsym} 
\usepackage{amssymb} 
\usepackage{amsmath}
\usepackage{amsxtra}
\usepackage[mathcal]{eucal} 
\usepackage{enumerate}
\usepackage{hyperref}
\usepackage{rotating}
\usepackage{caption}
\usepackage{lscape}
\usepackage{paralist} 
\usepackage{indentfirst}
\usepackage[dvipsnames,svgnames]{xcolor} 
\usepackage{setspace}
\usepackage{multicol} 
\usepackage{float}

\definecolor{codegreen}{rgb}{0,0.6,0}
\definecolor{codegray}{rgb}{0.5,0.5,0.5}
\definecolor{codepurple}{cmyk}{0.65, 1, 0, 0.35}
\definecolor{backcolour}{rgb}{0.95,0.95,0.92}

\usepackage{listings}
\usepackage{xcolor}

\lstset{
    language=R, % lenguaje de programación
    basicstyle=\small\ttfamily, % estilo básico
    commentstyle=\color{gray}, % estilo para los comentarios
    keywordstyle=\color{blue}, % estilo para las palabras clave
    stringstyle=\color{black}, % estilo para las cadenas de texto
    showstringspaces=false, % no mostrar espacios en cadenas de texto
    tabsize=4, % tamaño de la tabulación
    breaklines=true, % permitir el salto de línea automático
    postbreak=\mbox{\textcolor{red}{$\hookrightarrow$}\space}, % símbolo al final de una línea partida
    numbers=left, % mostrar números de línea
    numberstyle=\tiny\color{gray}, % estilo para los números de línea
    frame=single, % mostrar un borde alrededor del código
    backgroundcolor=\color{gray!10}, % color de fondo
}

\begin{document}

\begin{center}
    {\Large  Taller 2: Econometría }

Profesor:    Pedro Comte Hidalgo\\
Ayudante:     Barbara Buffo
\end{center}

\textbf{Instrucciones:} 

\begin{itemize}
    \item Usted dispone de 70 minutos  para intentar los 100 puntos del taller.
    \item Cuide la presentación y redacción de sus respuestas.
    \item Puede utilizar su computador y los apuntes tanto de clase como de ayudantía.
    \item Debe enviar un archivo en formato PDF y un R script (extensión .R)
\end{itemize}

\section{Trabajo Base de Datos (30 puntos)}

Para responder las preguntas siguientes, utilice la base de datos \textit{ceosal1} de la librería \textbf{wooldridge}. Asegúrese de cargarla correctamente (sin nulos) y de especificar el programa de forma rigurosa para evitar cambios por aleatoriedad.

\begin{enumerate}

\item (10 puntos) Examine la base y elabore una breve descripción de sus variables, distinguiendo entre dependientes, explicativas, continuas y \textit{dummies}. \textquestiondown Para qué tipo de estudio económico se emplea habitualmente esta base de datos?

\begin{lstlisting}
set.seed(1)
data(ceosal1)
ceosal1 <- na.omit(ceosal1)
?ceosal1
\end{lstlisting}


    \item (5 puntos) Genere una variable indicadora (\textit{dummy}) que tome el valor 1 cuando la variable \textit{sales} supere su promedio y 0 en caso contrario. Asigne a esta nueva variable un nombre que sea fácilmente identificable.
    

\begin{lstlisting}
ceosal1$sales_mayor_promedio <- ifelse(ceosal1$sales>mean(ceosal1$sales),1,0)
\end{lstlisting}

    \item (5 puntos) Cree una variable de la multuplicación entre la variable \textit{dummy} antes creada y la variable \textit{indus}. \textquestiondown Qué indica esta interacción?

\begin{lstlisting}
ceosal1$sales_mayor_promedio_indus <- ceosal1$sales_mayor_promedio*ceosal1$indus
\end{lstlisting}


    \item (5 puntos) Cree una variable de la multiplicación entre la variable continua \textit{roe}  y la variable \textit{dummy} antes creada. \textquestiondown Qué indica esta interacción?
\begin{lstlisting}
ceosal1$sales_mayor_promedio_roe <- ceosal1$sales_mayor_promedio*ceosal1$roe
\end{lstlisting}


    \item (5 puntos) Cree una variable que indique si el compañia corresponde a industria (\textit{indus}) ó si \textit{finance}.  Asigne a esta nueva variable un nombre que sea fácilmente identificable.
\begin{lstlisting}
ceosal1$indus_o_finance <- ceosal1$indus + ceosal1$finance
\end{lstlisting}


\end{enumerate}

\section{Modelos de Regresión Lineal (70 puntos)}

\begin{enumerate}

    \item (15 puntos) Genere un modelo de regresión lineal, sobre alguna de las variables dependientes, con respecto a las variables \textit{sales}, \textit{roe} y \textit{ros}. Interprete sus resultados y la significancia de los regresores. \textquestiondown El modelo tiene significancia global al 5\%?
\begin{lstlisting}
modelo <- lm(salary ~ sales+roe+ros, data=ceosal1)
summary(modelo)
\end{lstlisting}
\item (10 puntos) Al modelo de regresión lineal anterior, agregue las variables creadas en el inciso 1.2 y 1.4, interprete sus resultados. \textquestiondown Cómo afectan los regresores de estas variables en el salario?
\begin{lstlisting}
modelo <- lm(salary ~ sales+roe+ros+sales_mayor_promedio+ sales_mayor_promedio_roe, data=ceosal1)
summary(modelo)
\end{lstlisting}


\item (15 puntos) Realiza un test de hipotesis para evaluar si las dos variables que se agregaron son significativas para el modelo simultaneamente. Use significancia del 5\%.

\begin{lstlisting}
F_critico <- qf(0.95, 2 , 203)
R2_r <-0.03191
R2_ur <-0.04059
F_est <- (R2_ur-R2_r)/(1-R2_ur)*203/2
if (F_est>F_critico) {"Se rechaza H0"} else {"No se rechaza H0"}
\end{lstlisting}

\item (15 puntos) Volviendo a la regresión del salario en función de \textit{sales}, \textit{roe} y \textit{ros}. Considere dos grupos, Grupo 1: Aquellos clientes que \textit{sales} es mayor al promedio, y Grupo 2: Aquellos clientes que \textit{sales} es menor al promedio. Realice un test de hipotesis para grupos al 5\% de significancia. Use el metodo que más le acomode. 

\begin{lstlisting}
#haremos comparación de grupos mediante test F de significancia de variables
#para esto debemos crear variables que estan relacionadas con la dummy
ceosal1$sales_D <- ceosal1$sales*ceosal1$sales_mayor_promedio
ceosal1$roe_D <- ceosal1$roe*ceosal1$sales_mayor_promedio
ceosal1$ros_D <- ceosal1$ros*ceosal1$sales_mayor_promedio

modelo <- lm(salary ~ sales+roe+ros+sales_mayor_promedio+sales_D+ roe_D+ros_D, data=ceosal1)
summary(modelo)

F_critico <- qf(0.95, 4,201)
R2_ur <- 0.06028
R2_r <- 0.03191
F_est <- (R2_ur-R2_r)/(1-R2_ur)*201/4
if (F_est>F_critico) {"Se rechaza H0"} else {"No se rechaza H0"}
\end{lstlisting}

\item (15 puntos) Usando la regresión del salario en función de \textit{sales}, \textit{roe} y \textit{ros}, se distinguen los grupos mediante las \textit{dummies}; \textit{indus}, \textit{finance}, \textit{consprod} y \textit{utility}, que corresponden al tipo de empresa. Realice mediante test de Chow, con 5\% de significancia, si los modelos para cada grupo son iguales. 
\begin{lstlisting}
ceosal1_indus <- subset(ceosal1, indus==1)
ceosal1_finance <- subset(ceosal1, finance==1)
ceosal1_consprod <- subset(ceosal1, consprod==1)
ceosal1_utility <- subset(ceosal1, utility==1)

modelo <- lm(salary ~ sales+roe+ros, data=ceosal1)
modelo_indus <- lm(salary ~ sales+roe+ros, data=ceosal1_indus)
modelo_finance <- lm(salary ~ sales+roe+ros, data=ceosal1_finance)
modelo_consprod <- lm(salary ~ sales+roe+ros, data=ceosal1_consprod)
modelo_utility <- lm(salary ~ sales+roe+ros, data=ceosal1_utility)

SSR_r <- sum(residuals(modelo)^2)
SSR_indus <- sum(residuals(modelo_indus)^2)
SSR_finance <- sum(residuals(modelo_finance)^2)
SSR_consprod <- sum(residuals(modelo_consprod)^2)
SSR_utility <- sum(residuals(modelo_utility)^2)
SSR_ur <- SSR_indus+SSR_finance+SSR_consprod+SSR_utility

G <- 4
k <- 3
n <- nrow(ceosal1)
gl1 <- (G-1)*(k+1)
gl2 <- n - G*(k+1)

F_critico <- qf(0.95,gl1,gl2)
F_est <- (SSR_r - SSR_ur)/SSR_ur*gl2/gl1
if (F_est > F_critico) {"Se rechaza H0"} else {"No se rechaza H0"}
\end{lstlisting}


\end{enumerate}

\end{document}
